\documentclass{article}
\usepackage[T1]{fontenc}
\usepackage{lmodern}
\usepackage{graphicx}
\usepackage{amsmath,amssymb}
\usepackage[a4paper,margin=2cm]{geometry}
\usepackage[absolute,overlay]{textpos}
\setlength{\TPHorizModule}{1mm}
\setlength{\TPVertModule}{1mm}
\usepackage[indonesian]{babel}
\usepackage{hyperref}
\setlength{\emergencystretch}{2em}
% unit: Halaman 6

\begin{document}

% === Halaman 6 (llm) ===

\begin{itemize}
  \item Sebagai contoh, misalkan $p = 7$, maka $a = 3$ adalah akar primitif dari 7 karena
  \[ 
  \begin{aligned}
  3^1 \pmod 7 &= 3; & 3^2 \pmod 7 &= 2; & 3^3 \pmod 7 &= 6 \\
  3^4 \pmod 7 &= 4; & 3^5 \pmod 7 &= 5; & 3^6 \pmod 7 &= 1
  \end{aligned}
  \]
  \item Jadi, semua perpangkatan dari 3 menghasilkan nilai-nilai yang berbeda $(3, 2, 6, 4, 5, 1)$, semua bilangan di dalam modulus 7 terjadi satu kali.
  \item Perpangkatan berikutnya, $3^7 \pmod 7, 3^8 \pmod 7, \dots,$ akan kembali berulang menghasilkan nilai-nilai tersebut. Panjang satu siklus tidak lebih dari $7 - 1 = 6$.
  \item Secara umum, untuk $p$ bilangan prima, maka panjang siklus tidak lebih dari $p - 1$.
\end{itemize}

\end{document}
